\subsection{Factoring Out the GCF} 
The first step in factoring any polynomial is to find any monomial factor. We call this \term{factoring out the GCF}. Let's see an example, then summarize the process.
\begin{example}
Factor out the GCF from the polynomial $6x^5y^2-18x^4y^3+12y^4$
\begin{lpic}{}{6}
\end{lpic}
\end{example}
\begin{procedure}{Factoring Out the GCF}
\begin{enumerate*}
\item Find the GCF of the polynomial's terms.
\item Divide each term by the GCF to determine the remaining factor, using properties of exponents.
\item Write the original polynomial as the GCF multiplied by the polynomial made up of the remaining factors.
\item To check, note that:
\begin{enumerate*}
\item The remaining factors have no common factors.
\item When you multiply, you get the original polynomial.
\end{enumerate*}
\end{enumerate*}
\end{procedure}